% Generated draft: requires mathematical and source review.
\documentclass[11pt]{article}
\usepackage[margin=25mm]{geometry}
\usepackage{fontspec}
\setmainfont{Latin Modern Roman}
\usepackage{amsmath,amssymb,mathtools,unicode-math}
\setmathfont{Latin Modern Math}
\usepackage{enumitem,needspace,xurl,hyperref}
\setlength{\parindent}{0pt}
\setlength{\parskip}{5pt}
\setlength{\emergencystretch}{3em}
\newcommand{\N}{\mathbb{N}}
\newcommand{\Z}{\mathbb{Z}}
\newcommand{\Q}{\mathbb{Q}}
\newcommand{\R}{\mathbb{R}}
\newcommand{\C}{\mathbb{C}}
\newcommand{\F}{\mathbb{F}}
\newcommand{\A}{\mathbb{A}}
\newcommand{\PP}{\mathbb{P}}
\newcommand{\E}{\mathbb{E}}
\newcommand{\eps}{\varepsilon}
\newcommand{\dd}{\,\mathrm{d}}
\begin{document}
\section*{UnsolvedMath ID: 3700029}
\textbf{Source problem code:} catalogue entry
\subsection{Definitions and mathematical statement}
Let $D=\{z\in\C:|z|<1\}$. With $m$ interpreted as planar area measure, define, for integers $n\geq1$,
\[
\phi(n)=\sup_{\substack{p\in\C[z]\\\deg p=n}}
\int_D\frac{|p'(z)|}{1+|p(z)|^2}\,\mathrm{d}m(z),
\qquad
\alpha=\limsup_{n\to\infty}\frac{\log\phi(n)}{\log n}.
\]
Only $n\geq2$ enters the limiting quotient. Multiplying area measure by a fixed positive constant does not alter this exponent. The source imposes no monicity, coefficient, or boundary normalization on the degree-$n$ polynomials.

For a regular compact set $E$ in the source's potential-theoretic sense, the record refers to a constant $\beta_E$ defined by the growth of lengths of level curves of its Green function, and then sets
\[
\beta=\sup\{\beta_E:E\text{ is connected and is an admissible regular compact set}\}.
\]
The question asks for a connection between these analytic and geometric exponents, in particular whether $\alpha=\beta$.

The source does not give the formula for $\beta_E$, the Green-function normalization and domain, the limiting level parameter, the length convention, or the admissibility conditions on $E$. Thus the last two displays retain the source's named quantity but do not define a self-contained geometric exponent. These missing data must be recovered before the proposed equality can be treated as a precise independent statement.
\subsection{Short English statement}
Compare the polynomial spherical-derivative integral growth exponent $\alpha$ with the maximal Green-level-curve length exponent $\beta$ for connected compact sets. Is $\alpha=\beta$? The supplied record omits the definition of the geometric exponent.
\subsection{Sources}
\begin{itemize}
\item Local UnsolvedMath registry, permanent ID 3700029, exact source code catalogue entry, Littlewood constants, Question A; proposed by Alexandre Eremenko and cited proposers, proposed year 2002. Registry status partially solved, classification PARTIAL-PROGRESS.
\item Alexandre Eremenko, \emph{My favorite unsolved problems}, \url{https://www.math.purdue.edu/~eremenko/uns1.html}. Supplied source item: ``Littlewood constants, Question A''; extraction pdf-text, accessed July 29, 2026. The research summary names Eremenko's \emph{littlewood.pdf} but supplies no separate full citation or URL for it. No PDF was accessed here.
\item The supplied research summary describes established connections and bounds, and says the equality remains unresolved. It reports $\beta\leq1/4$ and $\alpha\leq3+2\sqrt2$, attributing results generally to Eremenko and predecessors without a precise theorem or reference. These numerical claims and their normalization require verification; they are not used to complete the missing definition of $\beta_E$. Its description of $\phi(n)$ as growing ``roughly like $n^\alpha$'' is not substituted for the exact limsup definition.
\item Provenance states the record was retained after the author-maintained index's explicit solved/counterexample updates through the access date. Current status and rights are flagged NEEDS\_REVIEW: public access was verified, but redistribution permission was not established.
\end{itemize}
\end{document}
